% !PRINTED 463
\weberpage{463}
Die drei Argumente von (2) sind also die Wurzeln von \"aquivalenten quadratischen Formen.

Wir unterscheiden die drei F\"alle:

1. $D\equiv 0\ (\mathrm{mod}\ 4)$. Hier ist $B$ gerade, und $A$ und $C$ k\"onnen nicht beide gerade sein; von den drei Formen (3), (5) hat also die eine zwei ungerade \"au{\ss}ere Koeffizienten, die beiden anderen haben einen geraden und einen ungeraden \"au{\ss}eren Koeffizienten.

2. Ist $D\equiv 1\ (\mathrm{mod}\ 8)$, so ist $B$ ungerade, $B^2\equiv 1\ (\mathrm{mod}\ 8)$, $AC\equiv 0\ (\mathrm{mod}\ 2)$, also ist wenigstens einer der beiden Koeffizienten gerade, und unter den drei Formen (3), (5) ist eine, aber auch nur eine, die zwei gerade \"au{\ss}ere Koeffizienten hat.

3. Ist $D\equiv 5\ (\mathrm{mod}\ 8)$, so ist $4AC\equiv 4\ (\mathrm{mod}\ 8)$ und die beiden \"au{\ss}eren Koeffizienten $A$ und $C$ sind \gesperrt{ungerade}.

Wenn es nun gelingt, eine ganzzahlige Gleichung aufzustellen, der von den drei Wurzeln (2) nur die eine gen\"ugt, so folgt daraus, da{\ss} diese \gesperrt{eine} Wurzel rational durch $j(\omega)$ ausdr\"uckbar und also eine Klasseninvariante ist.

Es besteht zwischen den Funktionen
\weq{6}{u = f(\omega)^{24}, \quad v = f\left(\frac{c+\partial\omega}{a}\right)^{24} \quad (c\equiv 0)\ (\mathrm{mod}\ 2)}
f\"ur einen ungeraden Transformationsgrad $n$ (\S\S\,73, 74) eine Transformationsgleichung:
\[
\Phi_n(u,v) = 0,
\]
und die Gleichung:
\weq{7}{\Phi_n(u,u) = 0}
ist also nach \S\,53 nur dann befriedigt, wenn f\"ur eine der Gr\"o{\ss}en (6):
\weq{8}{\frac{c+\partial\omega}{a} = \frac{\gamma+\delta\omega}{\alpha+\beta\omega},}
worin
\weq{9}{\begin{aligned}
&\begin{pmatrix} \alpha, & \beta \\ \gamma, & \delta \end{pmatrix} \equiv \begin{pmatrix} 1, & 0 \\ 0, & 1 \end{pmatrix} \quad (\mathrm{mod}\ 2)\\
\text{oder}\quad &\begin{pmatrix} \alpha, & \beta \\ \gamma, & \delta \end{pmatrix} \equiv \begin{pmatrix} 0, & 1 \\ -1, & 0 \end{pmatrix} \quad (\mathrm{mod}\ 2)
\end{aligned}}
eine zur ersten oder zur zweiten Klasse (\S\,36) geh\"orige lineare Transformation ist.

Die Vergleichung von (8) mit (3) f\"uhrt aber, wie oben, zu den Bedingungen:
%CONT>>
